% Part: sets-functions-relations
% Chapter: functions
% Section: isomorphisms

\documentclass[../../../include/open-logic-section]{subfiles}


\begin{document}

\olfileid[pt]{sfr}{fun}{iso}
\olsection{Isomorfismo}

\begin{explain}
Um \emph{isomorfismo} é uma bijeção que preserva a estrutura dos
conjuntos que relaciona, onde estrutura é uma questão das relações
que valem entre os !!{element}s dos conjuntos. Considere os
seguintes dois conjuntos $X=\{1,2,3\}$ e $Y=\{4,5,6\}$. Estes conjuntos são
ambos estruturados pelas relações sucessor, menor que e maior
que. Um isomorfismo entre os dois conjuntos é !!a{bijection} que
preserva essas estruturas. Assim, uma função !!{bijective} $f \colon X
\to Y$ é um isomorfismo se $i<j$ sse $f(i)<f(j)$, $i>j$ sse
$f(i)>f(j)$, e $j$ é o sucessor de $i$ sse $f(j)$ é o
sucessor de $f(i)$.
\end{explain}

\begin{defn}[Isomorfismo]
Seja $U$ o par $\langle X, R\rangle$ e $V$ o par $\langle
Y, S\rangle$ tais que $X$ e $Y$ são conjuntos e $R$ e $S$ são
relações sobre $X$ e $Y$ respectivamente. Uma !!{bijection} $f$ de $X$ para
$Y$ é um \emph{isomorfismo} de $U$ para $V$ sse preserva a
estrutura relacional, isto é, para quaisquer $x_{1}$ e $x_{2}$ em $X$,
$\tuple{x_1,x_2} \in R$ sse $\tuple{f(x_1),f(x_2)} \in S$.
\end{defn}

\begin{ex}
Considere os seguintes dois conjuntos $X=\{1,2,3\}$ e $Y=\{4,5,6\}$ e
as relações menor que e maior que. A função $f\colon X \to
Y$ onde $f(x) = 7-x$ é um isomorfismo entre $\tuple{X,<}$ e
$\tuple{Y,>}$.
\end{ex}

\end{document}
